\subsection{Connected fan pieces in a smaller closed square}
\label{sec:area_law_fan_closed_half}

The following local assertions concern the eight-triangle fan used in the
initial-star construction of \cite[Section~11]{OpenAI2026AreaLaw}.
The origin, dyadic exponent, cell index and binary coloring are arbitrary.
The exponent zero and the empty change set are included; distinct runs
carrying the same color remain distinct.

% Source: openai/math@adc7f1241b42e322a6451854ab7e4b4c146bf78a,
% 10-geometry.tex, geometry:initial-stars, lines 352--370, and
% prop:two-families, lines 308--323.

\begin{definition}[Retained pieces in the smaller closed square]
    \label{def:area_law_fan_closed_half_piece}
    \lean{TNLean.PEPS.AreaLaw.Geometry.cellFanClosedHalfPiece}
    \leanok
    \uses{def:area_law_fan_cut_piece}
    Let $o\in\R^2$, $\ell\in\N$, $z\in\Z^2$ and the colors
    $f_i\in\{0,1\}$ be arbitrary. Use the center $c$, half-side
    $r=2^\ell/2>0$, triangles $P_i$, change set $A$, deleted radials
    $L$ and retained pieces $S_i$ of
    Definition~\ref{def:area_law_fan_cut_piece}. Define
    \begin{align}
        K_i
         & =S_i\cap\overline B_\infty(c,r/2)
        =(P_i\setminus L)\cap\overline B_\infty(c,r/2).
    \end{align}
    The closed maximum-norm ball is the closed square of radius $r/2$.
    The equality follows from
    $\overline B_\infty(c,r/2)\subset B_\infty(c,r)$.
    When $A=\varnothing$, the center belongs to every $K_i$.
\end{definition}

\begin{theorem}[Intersection in the smaller closed square]
    \label{thm:area_law_fan_closed_half_intersection}
    \lean{TNLean.PEPS.AreaLaw.Geometry.cellFanPolygon_cut_inter_closedHalfBall_nonempty_of_adj}
    \leanok
    \uses{def:area_law_fan_closed_half_piece,def:area_law_fan_runs}
    If $i,j$ are consecutive perimeter positions with $f_i=f_j$, then
    $K_i\cap K_j$ is nonempty. Either $j=\tau(i)$
    or $i=\tau(j)$ is permitted, and the change set may be empty.
\end{theorem}
\begin{proof}
    \leanok
    \uses{def:area_law_fan_closed_half_piece,
        thm:area_law_cell_fan_successor_iff,
        thm:area_law_cell_fan_base_radius,
        thm:area_law_cell_fan_end_same_ray,
        thm:area_law_fan_noncentral_radial_incidence}
    Suppose first that $j=\tau(i)$, so $b_i=a_j$, and set
    $m=(c+b_i)/2$. The base-radius identity gives
    $\|m-c\|_\infty=r/2>0$, so $m\ne c$ and
    $m\in\overline B_\infty(c,r/2)\subset B_\infty(c,r)$.
    The noncentral radial-incidence characterization places $m$ in
    both $P_i$ and $P_j$.

    If $m\in[c,b_t]$, then $b_t-c$ and $b_i-c$ have the same
    directed ray. The distinctness of the endpoint directions gives
    $t=i$. But $f_i=f_j=f_{\tau(i)}$, so $i\notin A$.
    Thus $m\notin L$ and $m\in K_i\cap K_j$.
    For $i=\tau(j)$, interchange the indices and use symmetry of
    intersection.
\end{proof}

\begin{theorem}[Connectedness of a retained piece]
    \label{thm:area_law_fan_closed_half_piece_connected}
    \lean{TNLean.PEPS.AreaLaw.Geometry.cellFanClosedHalfPiece_isConnected}
    \leanok
    \uses{def:area_law_fan_closed_half_piece}
    Every $K_i$ is nonempty and connected.
\end{theorem}
\begin{proof}
    \leanok
    \uses{def:area_law_fan_closed_half_piece,
        thm:area_law_fan_cut_piece_connected,
        thm:area_law_fan_radial_union_deletion_convex}
    Apply Theorem~\ref{thm:area_law_fan_cut_piece_connected} at
    radius $r/2$ to obtain a point of
    $(P_i\setminus L)\cap B_\infty(c,r/2)$, which belongs to
    $K_i$. The deleted triangle is convex by
    Theorem~\ref{thm:area_law_fan_radial_union_deletion_convex}.
    Both centered balls are convex, so their intersection with the
    deleted triangle is a nonempty convex set and hence connected.
\end{proof}

\begin{theorem}[Connectedness of a retained run]
    \label{thm:area_law_fan_closed_half_run_connected}
    \lean{TNLean.PEPS.AreaLaw.Geometry.cellFanRunRegion_cut_inter_closedHalfBall_isConnected}
    \leanok
    \uses{def:area_law_fan_closed_half_piece,def:area_law_fan_runs}
    Retain an adjacency between consecutive triangles precisely when
    their colors agree. For any connected component $R$ of this graph,
    put $T_R=\bigcup_{i\in R}P_i$. Then
    $(T_R\setminus L)\cap\overline B_\infty(c,r/2)$ is nonempty
    and connected. Components carrying the same color remain distinct.
\end{theorem}
\begin{proof}
    \leanok
    \uses{def:area_law_fan_closed_half_piece,def:area_law_fan_runs,
        thm:area_law_fan_closed_half_piece_connected,
        thm:area_law_fan_closed_half_intersection}
    Each $K_i$ is nonempty and connected by
    Theorem~\ref{thm:area_law_fan_closed_half_piece_connected}, and
    Theorem~\ref{thm:area_law_fan_closed_half_intersection} gives
    a nonempty intersection along every retained edge.
    The component $R$ is nonempty, and any two of its vertices are
    joined by a finite path within $R$. Successive unions along these
    paths show that $\bigcup_{i\in R}K_i$ is nonempty and
    connected. Distributing deletion and intersection over the union
    gives
    \begin{align}
        (T_R\setminus L)\cap\overline B_\infty(c,r/2)
         & =\bigcup_{i\in R}K_i.
        \label{eq:area_law_closed_half_run}
    \end{align}
    Hence the set on the left of~(\ref{eq:area_law_closed_half_run})
    is nonempty and connected. If $A=\varnothing$, no radial is removed
    and the center is retained.
\end{proof}
